\documentclass[12pt,a4paper]{article}

%------------------------------------------------
% Packages
%------------------------------------------------
\usepackage[T1]{fontenc}
\usepackage{lmodern}
\usepackage[french]{babel}
\usepackage[margin=2cm]{geometry}
\usepackage{amsmath, amssymb}
\usepackage{booktabs}
\usepackage{array}
\usepackage{microtype}
\usepackage{graphicx}
\usepackage{hyperref}

%------------------------------------------------
% Informations du document
%------------------------------------------------
\title{\textbf{Deep Q-Network (DQN) : Étude comparative  \\ CartPole-v1, GridWorld-v0, LunarLander-v2}}
\author{Moubarak LASSISSI}
\date{Janvier 2026}

\begin{document}

\maketitle
\tableofcontents
\newpage

%================================================
\section{Introduction}

Le \textit{Deep Q-Network} (DQN) est un algorithme fondamental en apprentissage par renforcement profond. Il combine le Q-learning traditionnuel avec des réseaux de neurones profonds, permettant aux agents d'apprendre à partir d'observation brutes dans des espaces d'états continus avec des espaces d'actions discrets.

Ce rapport présente une implémentation complète et une analyse comparative du DQN appliqué à trois environnements de complexité croissante :
\begin{enumerate}
    \item \textbf{CartPole-v1} : équilibrage de perche (complexité basse),
    \item \textbf{GridWorld-v0} : navigation en grille 2D (complexité moyenne),
    \item \textbf{LunarLander-v2} : atterrissage lunaire (complexité haute).
\end{enumerate}

\subsection{Motivation et objectifs}

L'efficacité du DQN repose sur plusieurs composants clés. Comprendre l'importance de chacun est essentiel pour optimiser les performances :

\begin{itemize}
    \item \textbf{Experience Replay :} le stockage et l'échantillonnage aléatoire des transitions passées réduit la corrélation entre les samples d'apprentissage, améliorant la stabilité et l'efficacité de l'algorithme.

    \item \textbf{Target Network :} l'utilisation d'un réseau secondaire pour calculer les valeurs cible réduit les oscillations et stabilise significativement l'apprentissage.

    \item \textbf{Double DQN :} réduit le biais de surestimation des Q-values en découplant la sélection et l'évaluation des actions.

    \item \textbf{Prioritized Replay (PER) :} accélère l'apprentissage en échantillonnant prioritairement les transitions avec erreur TD élevée.
\end{itemize}

\subsection{Approche méthodologique}

L'objectif principal de ce rapport est de conduct une \textit{ablation study} systématique mesurant l'impact isolé et combiné de chaque composant du DQN. Nous évaluons \textbf{5 variantes} :

\begin{enumerate}
    \item \textbf{Baseline :} Q-learning simple (aucun optimisation).
    \item \textbf{ReplayOnly :} + Experience Replay seul.
    \item \textbf{TargetOnly :} + Target Network seul.
    \item \textbf{Standard :} Experience Replay + Target Network (DQN classique).
    \item \textbf{Advanced :} Double DQN + Prioritized Replay + Target Network.
\end{enumerate}

Cette étude comparative nous permettra d'identifier les contributions spécifiques de chaque technique et de déterminer les conditions où chacune apporte un bénéfice significatif de performance.

%================================================
\section{Algorithme DQN}

\subsection{Q-learning}

Le Q-learning vise à estimer la fonction valeur-action :

\[
Q(s,a) = \mathbb{E}\left[r + \gamma \max_{a'} Q(s',a')\right]
\]

avec la règle de mise à jour :

\[
Q(s,a) \leftarrow Q(s,a) + \alpha \left(r + \gamma \max_{a'} Q(s',a') - Q(s,a)\right)
\]

où $\alpha \in (0,1]$ est le taux d'apprentissage.

\subsection{Extension aux réseaux de neurones}

Dans le DQN, la fonction $Q$ est approximée par un réseau de neurones $\hat{Q}(s,a;\theta)$. La fonction de perte minimisée est :

\[
\mathcal{L}(\theta) =
\mathbb{E}\left[
\left(
r + \gamma \max_{a'} \hat{Q}(s',a';\theta^-)
- \hat{Q}(s,a;\theta)
\right)^2
\right]
\]

où $\theta^-$ correspond aux paramètres du \textit{target network}.

\subsection{Experience Replay}

Les transitions $(s,a,r,s')$ sont stockées dans un buffer de replay. À chaque étape d'apprentissage, un mini-batch est échantillonné aléatoirement :
\begin{itemize}
    \item réduit la corrélation entre les échantillons,
    \item améliore la stabilité de l'apprentissage,
    \item réutilise efficacement les expériences passées.
\end{itemize}

\subsection{Target Network}

Un second réseau, appelé \textit{target network}, calcule les cibles. Ses paramètres $\theta^-$ sont mis à jour périodiquement à partir du réseau principal.

\subsection{Double DQN}

Le Double DQN réduit la surestimation des Q-values. La cible devient :

\[
r + \gamma \hat{Q}(s', \arg\max_{a'} \hat{Q}(s',a';\theta); \theta^-)
\]

\subsection{Prioritized Experience Replay}

Plutôt que d'échantillonner uniformément, les transitions sont échantillonnées selon leur erreur TD :

\[
p_i \propto |\delta_i|^{\alpha}, \quad \delta_i = r + \gamma \max_{a'} \hat{Q}(s',a';\theta^-) - \hat{Q}(s,a;\theta)
\]

\subsection{Politique d'exploration}

L'agent utilise une politique $\varepsilon$-gloutonne :

\[
a_t =
\begin{cases}
\text{action aléatoire} & \text{avec probabilité } \varepsilon, \\
\arg\max_a \hat{Q}(s,a;\theta) & \text{sinon}.
\end{cases}
\]

Le paramètre $\varepsilon$ décroît progressivement.

%================================================
\section{Environnements}

%================================================
\subsection{CartPole-v1}

\subsubsection{Description}

CartPole-v1 est un problème classique de contrôle avec espace d'état continu et espace d'actions discret.

\begin{itemize}
    \item \textbf{État :} vecteur réel de dimension 4 :
    \begin{enumerate}
        \item position du chariot ($x$),
        \item vitesse du chariot ($\dot{x}$),
        \item angle de la perche ($\theta$),
        \item vitesse angulaire de la perche ($\dot{\theta}$).
    \end{enumerate}

    \item \textbf{Actions :} 2 actions discrètes :
    \begin{enumerate}
        \item pousser le chariot vers la gauche,
        \item pousser le chariot vers la droite.
    \end{enumerate}

    \item \textbf{Récompense :} +1 à chaque pas de temps dans l'épisode.

    \item \textbf{Terminalité :}
    \begin{itemize}
        \item perche dépasse 15°,
        \item chariot sort de l'écran ($|x| > 2.4$),
        \item 500 pas atteints (version v1).
    \end{itemize}

    \item \textbf{Complexité :} Basse ($\sim 10$ min apprentissage).
\end{itemize}

\subsubsection{Objectif}

L'agent cherche à apprendre une politique optimale $\pi^*$ maximisant :
\[
J(\pi) = \mathbb{E}\left[\sum_{t=0}^{T} \gamma^t r_t\right]
\]

Cible : reward $\geq 400$ sur 100 épisodes de test.

\subsubsection{Architecture réseau CartPole}

\[
\begin{array}{c}
\text{Input (4 features)} \\
\downarrow \\
\text{Linear}(4 \to 64) + \tanh \\
\downarrow \\
\text{Linear}(64 \to 64) + \tanh \\
\downarrow \\
\text{Linear}(64 \to 2) \quad \text{[Q-values]}
\end{array}
\]

%================================================
\subsection{GridWorld-v0}

\subsubsection{Description}

GridWorld est un problème de navigation en grille 2D avec observations spatiales.

\begin{itemize}
    \item \textbf{État :} grille $n \times m$ encodant :
    \begin{enumerate}
        \item position de l'agent,
        \item positions cibles (récompense),
        \item positions obstacles/murs.
    \end{enumerate}
    Feature extraction : $\text{MapFromDumpExtractor2}$ → 3 canaux $\times$ map size

    \item \textbf{Actions :} 4 actions discrètes :
    \begin{enumerate}
        \item haut, bas, gauche, droite.
    \end{enumerate}

    \item \textbf{Récompense :}
    \begin{itemize}
        \item +1 atteindre objectif,
        \item -0.001 par pas (pénalité de temps),
        \item -1 atteindre obstacle.
    \end{itemize}

    \item \textbf{Terminalité :}
    \begin{itemize}
        \item objectif atteint,
        \item 100 pas max (artificiel en training).
    \end{itemize}

    \item \textbf{Complexité :} Moyenne ($\sim 20$ min apprentissage).
\end{itemize}

\subsubsection{Architecture réseau GridWorld}

\[
\begin{array}{c}
\text{Input (extracted features)} \\
\downarrow \\
\text{Linear}(n_{\text{feat}} \to 30) + \tanh \\
\downarrow \\
\text{Linear}(30 \to 30) + \tanh \\
\downarrow \\
\text{Linear}(30 \to 4) \quad \text{[Q-values]}
\end{array}
\]

%================================================
\subsection{LunarLander-v2}

\subsubsection{Description}

LunarLander est un problème de contrôle continu : atterrir un module lunaire sans crash.

\begin{itemize}
    \item \textbf{État :} vecteur 8D :
    \begin{enumerate}
        \item position $(x, y)$,
        \item vitesse $(\dot{x}, \dot{y})$,
        \item angle $\theta$, vitesse angulaire $\dot{\theta}$,
        \item indicateurs moteur gauche/droit.
    \end{enumerate}

    \item \textbf{Actions :} 4 actions discrètes :
    \begin{enumerate}
        \item rien, moteur droit, moteur gauche, moteur central.
    \end{enumerate}

    \item \textbf{Récompense :}
    \begin{itemize}
        \item jusqu'à 275 pour atterrir près de la cible,
        \item -300 crash,
        \item -0.3 par seconde combustible.
    \end{itemize}

    \item \textbf{Terminalité :}
    \begin{itemize}
        \item atterrissage (success/crash),
        \item 1000 pas max.
    \end{itemize}

    \item \textbf{Complexité :} Haute ($\sim 30$ min apprentissage).
\end{itemize}

\subsubsection{Architecture réseau LunarLander}

\[
\begin{array}{c}
\text{Input (8 state variables)} \\
\downarrow \\
\text{Linear}(8 \to 200) + \tanh \\
\downarrow \\
\text{Linear}(200 \to 4) \quad \text{[Q-values]}
\end{array}
\]

%================================================
\section{Ablation Study}

L'objectif est de mesurer l'impact de chaque composant du DQN sur les trois environnements.

\subsection{Variantes testées}

\begin{enumerate}
    \item \textbf{Baseline :} Q-learning simple (pas de replay buffer, pas de target network).
    \item \textbf{ReplayOnly :} + Experience Replay (uniform sampling).
    \item \textbf{TargetOnly :} + Target Network (pas de PER).
    \item \textbf{Standard :} Experience Replay + Target Network (DQN classique).
    \item \textbf{Advanced (BONUS)} : Double DQN + Prioritized Replay + Target Network.
\end{enumerate}

\subsection{Protocole expérimental}

\begin{itemize}
    \item \textbf{Seed :} fixé à 5 pour reproducibilité.
    \item \textbf{Nombre d'épisodes :} 1,000,000 par variant/environnement.
    \item \textbf{Test frequency :} tous les 100-1000 épisodes selon l'environnement.
    \item \textbf{Test samples :} 5 épisodes par évaluation (moyennes stables).
    \item \textbf{Hyperparamètres constants :}
    \begin{itemize}
        \item Learning rate : 0.0001
        \item Discount factor : $\gamma = 0.99$
        \item Initial exploration : $\varepsilon_0 = 0.1$
        \item Batch size : 100
    \end{itemize}
\end{itemize}

\subsection{Métriques}

Pour chaque variante, on extrait :
\begin{itemize}
    \item \textbf{reward :} récompense cumulée par épisode (train).
    \item \textbf{rewardTest :} performance de test (stabilité).
    \item \textbf{qvalue\_mean :} estimation Q moyenne (stabilité réseau).
    \item \textbf{epsilon :} taux d'exploration (decay).
    \item \textbf{tdError\_mean :} erreur TD (convergence).
\end{itemize}

%================================================
\section{Résultats}

\subsection{CartPole-v1}

\subsubsection{Convergence par variante}

\textit{[Insérer graphique : Reward vs Episode pour les 5 variants]}

\begin{itemize}
    \item \textbf{Baseline :} convergence lente, instabilité.
    \item \textbf{ReplayOnly :} amélioration significative.
    \item \textbf{TargetOnly :} stabilité accrue.
    \item \textbf{Standard :} convergence rapide et stable.
    \item \textbf{Advanced :} performance optimale.
\end{itemize}

\subsubsection{Performance de test}

\textit{[Insérer graphique : RewardTest vs Episode]}

Tableau récapitulatif :

\begin{center}
\begin{tabular}{|l|r|r|r|}
\hline
\textbf{Variante} & \textbf{Reward moyen} & \textbf{Reward max} & \textbf{Stabilité} \\
\hline
Baseline & XXX & XXX & $\pm$ XXX \\
ReplayOnly & XXX & XXX & $\pm$ XXX \\
TargetOnly & XXX & XXX & $\pm$ XXX \\
Standard & XXX & XXX & $\pm$ XXX \\
Advanced & XXX & XXX & $\pm$ XXX \\
\hline
\end{tabular}
\end{center}

%------------------------------------------------
\subsection{GridWorld-v0}

\subsubsection{Convergence par variante}

\textit{[Insérer graphique : Reward vs Episode pour les 5 variants]}

Observations :
\begin{itemize}
    \item Impact du Target Network plus marqué que sur CartPole.
    \item PER apporte un gain significatif en rapidité.
\end{itemize}

\subsubsection{Tableau comparatif}

\begin{center}
\begin{tabular}{|l|r|r|r|}
\hline
\textbf{Variante} & \textbf{Reward moyen} & \textbf{Reward max} & \textbf{Stabilité} \\
\hline
Baseline & XXX & XXX & $\pm$ XXX \\
ReplayOnly & XXX & XXX & $\pm$ XXX \\
TargetOnly & XXX & XXX & $\pm$ XXX \\
Standard & XXX & XXX & $\pm$ XXX \\
Advanced & XXX & XXX & $\pm$ XXX \\
\hline
\end{tabular}
\end{center}

%------------------------------------------------
\subsection{LunarLander-v2}

\subsubsection{Convergence par variante}

\textit{[Insérer graphique : Reward vs Episode pour les 5 variants]}

Observations :
\begin{itemize}
    \item Tâche complexe nécessite tous les composants.
    \item Double DQN critique pour éviter surestimation.
    \item PER essentiel pour convergence.
\end{itemize}

\subsubsection{Tableau comparatif}

\begin{center}
\begin{tabular}{|l|r|r|r|}
\hline
\textbf{Variante} & \textbf{Reward moyen} & \textbf{Reward max} & \textbf{Stabilité} \\
\hline
Baseline & XXX & XXX & $\pm$ XXX \\
ReplayOnly & XXX & XXX & $\pm$ XXX \\
TargetOnly & XXX & XXX & $\pm$ XXX \\
Standard & XXX & XXX & $\pm$ XXX \\
Advanced & XXX & XXX & $\pm$ XXX \\
\hline
\end{tabular}
\end{center}

%================================================
\section{Analyse comparative}

\subsection{Impact d'Experience Replay}

En comparant \textit{Baseline} vs \textit{ReplayOnly} :

\[
\text{Gain} = \frac{\text{Reward}_{ReplayOnly} - \text{Reward}_{Baseline}}{\text{Reward}_{Baseline}} \times 100\%
\]

\begin{center}
\begin{tabular}{|l|r|}
\hline
\textbf{Environnement} & \textbf{Gain moyen} \\
\hline
CartPole & XXX\% \\
GridWorld & XXX\% \\
LunarLander & XXX\% \\
\hline
\end{tabular}
\end{center}

\subsection{Impact du Target Network}

Comparaison \textit{Baseline} vs \textit{TargetOnly}.

\subsection{Synérgie Replay + Target}

Comparaison \textit{Standard} vs (\textit{ReplayOnly} + \textit{TargetOnly}).

\subsection{Efficacité du Double DQN + PER}

Comparaison \textit{Advanced} vs \textit{Standard}.

%================================================
\section{Conclusion}

\subsection{Résumé des résultats}

\begin{itemize}
    \item \textit{CartPole} : problème simple, tous les variants convergent.
    \item \textit{GridWorld} : Target Network crucial, PER apporte stabilité.
    \item \textit{LunarLander} : Double DQN + PER essentiels, baseline échoue.
\end{itemize}

\subsection{Points clés}

\begin{enumerate}
    \item Experience Replay seul n'est pas suffisant sur tâches complexes.
    \item Target Network stabilise significativement l'apprentissage.
    \item Double DQN réduit surestimation critique pour tâches longues.
    \item Prioritized Replay accélère convergence sur problèmes difficiles.
\end{enumerate}

\subsection{Perspectives}

\begin{itemize}
    \item Implémenter Dueling DQN pour décomposer $V(s)$ et $A(s,a)$.
    \item Tester sur Atari (images) avec convolutions.
    \item Analyser variance and bias du réseau de neurones.
\end{itemize}

\end{document}
